Als
\definitionsverweis {Einheitsnormalenfeld}{}{}
können wir vom Gradienten der beschreibenden Funktion, also von

\mavergleichskettedisp
{\vergleichskette
{ G(x,y)
}
{ =} { \begin{pmatrix}  2 \alpha x  \\ 2 \beta y    \end{pmatrix}
}
{ } {
}
{ } {
}
{ } {
}
}
{}{}{}
ausgehen, durch Normalisierung erhalten wir

\mavergleichskettedisp
{\vergleichskette
{ N(x,y)
}
{ =} {  { \frac{   1 }{  \sqrt{ \alpha^2 x^2 + \beta^2 y^2 }  } } \begin{pmatrix}   \alpha x  \\  \beta y    \end{pmatrix}
}
{ } {
}
{ } {
}
{ } {
}
}
{}{}{.}
Die Gauß-Abbildung ist also
\maabbeledisp {\varphi} {Y = \left\{  \begin{pmatrix}  x  \\ y   \end{pmatrix}   \in \R^2  \mid   \alpha x^2+\beta y^2 = 1        \right\}  } { S^1
} { \begin{pmatrix}  x  \\ y   \end{pmatrix} } { { \frac{   1 }{  \sqrt{ \alpha^2 x^2 + \beta^2 y^2 }  } } \begin{pmatrix}   \alpha x  \\  \beta y    \end{pmatrix}
} {.} Die Surjektivität der Abbildung folgt aus
[[Hyperfläche/Glatt und kompakt/Gauß-Abbildung surjektiv/Fakt|Fakt]].
Zum Nachweis der Injektivität betrachten wir die Bedingung

\mavergleichskettedisp
{\vergleichskette
{ { \frac{   1  }{  \sqrt{ \alpha^2 x^2 + \beta^2 y^2 }  } } \begin{pmatrix}   \alpha x  \\  \beta y    \end{pmatrix}
}
{ =} { { \frac{   1  }{  \sqrt{ \alpha^2 z^2 + \beta^2 w^2 }  } } \begin{pmatrix}   \alpha z  \\  \beta w    \end{pmatrix}
}
{ } {
}
{ } {
}
{ } {
}
}
{}{}{.}
Das bedeutet, dass
\mathkor {} {\begin{pmatrix}   \alpha x  \\  \beta y    \end{pmatrix}} {und} {\begin{pmatrix}   \alpha z  \\  \beta w    \end{pmatrix}} {}
linear abhängig sind mit einem positiven Faktor $c$ und dass ihre Normierungen gleich sind. Dann ist aber

\mavergleichskettedisp
{\vergleichskette
{ \begin{pmatrix}   \alpha x  \\  \beta y    \end{pmatrix}
}
{ =} { \begin{pmatrix}   \alpha z  \\  \beta w    \end{pmatrix}
}
{ } {
}
{ } {
}
{ } {
}
}
{}{}{}
und daher

\mavergleichskette
{\vergleichskette
{ x
}
{ = }{ z
}
{  }{
}
{    }{
}
{   }{
}
}
{}{}{}
und

\mavergleichskette
{\vergleichskette
{ y
}
{ = }{ w
}
{  }{
}
{    }{
}
{   }{
}
}
{}{}{.}

 [[Kategorie:Latexseite]]
